% =====================================================================
%  Security posture: the threat model, what the code enforces, and what
%  a deployment would still have to do.
%  SCOPE: complements Sec. \ref{sec:phaseC_security}, which describes the
%         cryptographic construction. This one is about everything the
%         construction does not by itself achieve.
% =====================================================================
\section{Security Posture}
\label{sec:security_posture}

\subsection{Why this section exists}

Section~\ref{sec:phaseC_security} describes a sealed, signed report and
a signed request. Read on its own, that reads like a solved problem, and
it is not one. Encryption answers a narrow question --- can a third
party read this message, and did the named sender write it --- and a
deployed system fails in ways that leave both answers ``yes''.

An order captured off the broker and published again a week later
carries the Cloud's own signature. A stranger answering \emph{hold} to
every handover fills the robot's outbox and produces a campaign with no
data and no errors. A filed reading edited on disk was verified on
arrival and is verified by nothing afterwards. None of those are defeats
of the cryptography; all three were open in the version of this system
described by the previous section.

This section states the threat model, lists what the delivered code now
enforces, and --- at least as importantly --- says what it does not, so
that the boundary is a decision on record rather than something a reader
has to infer.

\subsection{Threat model}

The system is assumed to run on a small network the operator does not
exclusively control: a farm LAN with a shared wireless segment, other
machines on it, and physical access to the greenhouse by people who are
not the operator.

\begin{table}[!t]
\caption{Who is assumed capable of what. The middle row is the one that
drives most of the design: an attacker who never breaks a key can still
do a great deal.}
\label{tab:threats}
\centering\footnotesize
\begin{tabular}{@{}p{0.26\columnwidth}p{0.64\columnwidth}@{}}
\toprule
actor & assumed capability \\
\midrule
network observer & reads every packet on the segment; learns topics,
  timings and message sizes \\
network participant & connects to the broker and publishes on any topic;
  captures and replays messages verbatim \\
disk access & reads and writes files on the Cloud machine, including the
  store and the key directory \\
key holder & has one long-term private key, e.g.\ from a robot that was
  stolen or decommissioned \\
\bottomrule
\end{tabular}
\end{table}

Explicitly \emph{out} of scope: an attacker who can break P-256 or
AES-256-GCM, and one with root on the Cloud while it is running. Neither
is defensible at this scale, and pretending otherwise would make
everything else here less believable.

\subsection{What the delivered code enforces}

\paragraph{A signed order stops being an order}
\filename{agri/replay.py}. A signature says who wrote a message, not
that this is the first time it has been read. Requests now carry a
freshness window (\SI{120}{s}) \emph{and} a seen-set of request
identifiers. Neither works alone: freshness leaves the window open, and
a seen-set alone would have to remember every identifier forever and
across reboots --- which is exactly when a replay would be attempted.
Together they need no persistence and leave no gap, because an
identifier older than the window is refused on age before the set is
consulted. \param{request\_id} serves as the nonce; it is already
random, already unique and already inside the signature.

The failure mode this introduces is a clock, not an attacker, and the
code treats it that way: a future-dated order is reported as clock skew
and never as an attack, an accepted order already consuming half the
window logs a warning, and every refusal prints the disagreement in
seconds alongside \texttt{timedatectl}.

\paragraph{The handshake is signed in both directions}
The negotiation messages were left plain, on the reasoning that a forged
\emph{send} only makes the robot offer a reading it was about to offer
anyway. That reasoning missed the other grant. \emph{Hold} means keep
it, and the robot obeys by putting the reading in a bounded outbox.
Answered to every offer, it fills 48 slots, starts dropping the oldest
readings, and leaves a robot reporting itself perfectly healthy while
the campaign produces nothing --- for the cost of one MQTT publish.
Queries are now signed by the Cloud and replies by the node, each with a
key the other already holds.

\paragraph{The transport can be closed}
\filename{deploy/}. ECIES hides a reading; it does not hide that
\texttt{agri/v1/report/youbot-01} published \SI{6}{kB} at 14:02, and for
a commercial greenhouse that traffic pattern is a map of the operation.
Both programs accept a CA, a client certificate and a key; naming a CA
turns TLS on and moves the port from 1883 to 8883 without a second
decision. Neither side ever falls back to plaintext --- a
misconfiguration is fatal, because a system that silently downgrades is
worse than one with no TLS, since somebody will believe it. The
reference \filename{mosquitto.conf} does not listen on 1883 at all, and
the ACL's \texttt{\%u} patterns stop one node publishing as another and
stop any node issuing orders.

There is deliberately no password flag: \filename{/proc/<pid>/cmdline}
is world-readable, so a password in \texttt{argv} is one that
\texttt{ps} prints to anyone with a shell.

\paragraph{The dashboard credential leaves the URL}
A token in a query string lives in the address bar, the browser history,
the \texttt{Referer} of every outbound link, and any proxy log in
between --- for a credential that authorises driving the robot and
shutting the simulation down. It now moves into an
\texttt{HttpOnly; SameSite=Strict} cookie on first contact and the page
redirects to a clean address. Repeated wrong tokens earn a bounded
lockout. Disabling authentication is still possible and is now permitted
only when the server is bound to loopback: a property the program
checks, rather than a warning that scrolls past.

\paragraph{A filed reading that changes says so}
\filename{agri/cloud/store.py}. Every report is verified on arrival and
was then written in plain text, after which the cryptography had nothing
further to say about it. Each line now carries the SHA-256 of its own
content chained to the hash before it, so editing one reading breaks
that line, repairing that line breaks the next, and deleting a line is
caught by the chain even though every remaining line still hashes
correctly.

\paragraph{A key can be named and retired}
\filename{agri/trust.py}. Keys get a 64-bit fingerprint in four readable
groups, printed at startup by both programs, so ``do we hold the same
key'' is settled in ten seconds rather than by an afternoon of signature
errors. A revocation list is consulted at startup and both programs
\emph{refuse to run} on a listed key. \texttt{--rotate} archives the old
pair before touching anything --- reports already sealed to a private
key can be opened with nothing else --- and revokes the old public key
before writing the new one.

\subsection{What it does not do, and what that would cost}

\begin{table*}[!t]
\caption{Deliberate boundaries. Each row is a real weakness of the
delivered system; the third column is what closing it would take, so
that the boundary is a decision rather than an oversight.}
\label{tab:notdone}
\centering\footnotesize
\renewcommand{\arraystretch}{1.25}
\begin{tabular}{@{}p{0.20\textwidth}p{0.42\textwidth}p{0.30\textwidth}@{}}
\toprule
gap & why it is open & what would close it \\
\midrule
The revocation list is not distributed &
Revoking a key on the Cloud tells the robot nothing. Each machine holds
its own list and they are reconciled by \texttt{scp}. &
X.509 certificates with a CRL or OCSP, i.e.\ a real PKI. At two nodes
that is more machinery than the problem. \\

Key age is a reminder, not a validity period &
A raw PEM carries no expiry --- that is precisely what a certificate
adds. The birthday sits in an unsigned sidecar the key holder could
edit. &
The same PKI. Until then the honest claim is that an old key is
mentioned at every startup. \\

The hash chain lives on the same disk as the data &
An attacker with write access can recompute the whole chain. What the
chain buys is that a surgical edit is detectable and a forgery must
rewrite the file to its end. &
Pinning the tip somewhere out of reach: an off-box nightly copy, or a
signature from a key that is not on the machine.
\texttt{--verify} prints the tip so that this is possible. \\

The dashboard is plain HTTP &
The cookie is \texttt{HttpOnly} and \texttt{SameSite}, but not
\texttt{Secure} --- a \texttt{Secure} cookie on an HTTP origin is
silently discarded, which would look like authentication that does not
stick. &
A TLS reverse proxy, and then setting \texttt{Secure} and HSTS. No code
change. \\

One shared token, no accounts &
Every operator uses the same credential; there is no record of who
issued which request. \texttt{requests.csv} attributes orders to the
Cloud, not to a person. &
Per-operator credentials and an actor field inside the signed request.
The protocol has room; the deployment has no identity provider. \\

Dependencies are unpinned &
\texttt{pip install -e .} resolves \texttt{paho}, \texttt{cryptography},
\texttt{opencv} and \texttt{Pillow} at install time, so two machines can
run different code. &
A lock file, and a hash-checking install. Worth doing before this is
handed to anyone. \\

Acks and status are unsigned &
They only report, so forging one lies to a dashboard rather than
steering a robot --- a deliberate line, stated in
\filename{agri/protocol.py}. &
Signing them too. The cost is one verification per telemetry message at
\SI{0.2}{Hz}, which is nothing; the reason not to is that it would blur
the rule that makes the protocol legible. \\
\bottomrule
\end{tabular}
\end{table*}

\subsection{The rule the protocol now follows}

One line makes the message table readable without consulting the code:
\emph{everything that can change what the other side does is signed;
everything that only reports is not}. A request, a query and a reply all
decide something and all carry signatures. An acknowledgement and a
status are observations, and do not.

That rule is worth stating because the alternative --- signing
everything --- is the choice that looks safer and teaches nothing. It
would leave a reader unable to tell which messages actually matter, and
it is the same reflex that would have encrypted a request whose contents
are ``measure \station{P2,5R}''.

\subsection{What was found by writing this down}

Two of the gaps above were closed only because the exercise of
enumerating them made them impossible to keep deferring, and one defect
was found by a linter rather than a test: a dictionary literal in
\filename{agri/envelope.py} named two different values \texttt{plant},
which Python resolves last-wins in silence. Every report ever filed
carried a coordinate pair where the plant index belonged, beside a row
index that was a plain integer. Nothing downstream read the field, so
nothing ever broke --- it was simply wrong in the file, which is the
class of defect only found by reading the file or by a tool that reads
it for you.
